\section{Experiments}
\label{sec:experiments}

We first examine whether a fixed top-$L$ cutoff changes the measured effect of
query expansion. We then evaluate whether DESA improves retrieval effectiveness while
reducing the ranking depth required from dense and sparse retrieval. The
remaining experiments separate the contributions of the two channel operators
and test robustness across generators, dense encoders, reference counts, and
corpus sizes.

\subsection{Evaluation Protocol}

To measure cutoff sensitivity, we vary
$L\in\{10,20,50,100,200,500,1000\}$. At each cutoff, the top-$L$ dense and
sparse rankings are fused using equal-weight RRF
\citep{cormack-etal-2009-rrf}. We compute nDCG@10
\citep{jarvelin-kekalainen-2002-cumulated} and Recall@20, compare the ordering
of query constructions across cutoff depths, and measure whether each fixed-$L$
result reproduces the ordered top-20 from complete-list fusion.

We also examine whether the cutoff changes the judgment of an expansion
relative to the unexpanded query. For each query and generation draw, the
metric difference between an expansion and Original is labeled as a gain, tie,
or loss under both fixed-$L$ and complete-list fusion. We report how often this
label changes and separately count strict reversals, where a gain becomes a
loss or vice versa.

All remaining experiments use the complete-list fusion defined in
Section~\ref{sec:quality-access}. Retrieval effectiveness is measured by nDCG@10 and
Recall@20. Channel access is measured by the dense and sparse replay stopping
depths $(L_D,L_S)$. We report their changes from Original and the proportion of
queries for which both depths decrease. Sparse support size and exhaustion rate
are used to explain changes in lexical access.

\subsection{Datasets and Retrieval Setup}

We evaluate on seven English BEIR datasets \citep{thakur-etal-2021-beir}: SciFact
\citep{wadden-etal-2020-fact}, NFCorpus \citep{boteva-etal-2016-fulltext},
TREC-COVID \citep{voorhees-etal-2020-trec-covid}, FiQA
\citep{maia-etal-2018-fiqa}, ArguAna \citep{wachsmuth-etal-2018-retrieval},
Touch\'e-2020 \citep{bondarenko-etal-2020-touche}, and SCIDOCS
\citep{cohan-etal-2020-specter}. These datasets cover scientific claim
verification, biomedical and financial retrieval, argument retrieval, and
scientific document recommendation. Unless otherwise stated, the cutoff,
controlled, and mechanism analyses use all seven datasets, with each dataset
receiving equal weight in macro averages.

For the scale analysis, we construct nested TREC-COVID corpora containing
25,000, 50,000, 100,000, and all documents. Every snapshot retains the judged
relevant documents, while the remaining documents are added in a fixed order.
Each smaller corpus is therefore a strict subset of the larger ones.

The primary generator is Qwen2.5-7B-Instruct
\citep{qwen-team-2024-qwen25}. For each query, we perform three generation
draws and produce five complementary reference passages per draw. The generator
robustness experiment uses Mistral-7B-Instruct-v0.3
\citep{jiang-etal-2023-mistral} on up to 100 queries selected by a fixed hash
from FiQA, ArguAna, Touch\'e-2020, and SCIDOCS.

Dense retrieval uses BGE-small-en-v1.5 \citep{xiao-etal-2024-cpack}, with
normalized query and document vectors ranked by dot product. Sparse retrieval
uses Lucene BM25 through Pyserini 2.3.0, with $k_1=0.9$ and $b=0.4$
\citep{robertson-zaragoza-2009-bm25,lin-etal-2021-pyserini}. Contriever
\citep{izacard-etal-2022-contriever} is used as a second dense encoder. Score
ties are broken by document identifier.

\subsection{Comparison Methods}

The controlled comparison includes Original and Shared expansion. Original
submits the unexpanded query to both retrieval channels. Shared expansion uses
the same generated references as DESA but applies them directly to both
channels: dense retrieval averages the representations of the five query and
reference pairs, while sparse retrieval concatenates the original query and all
references. The comparison between Shared expansion and DESA therefore tests
whether dense and sparse retrieval benefit from different integration rules.

We also compare with three representative generative expansion methods. HyDE
encodes generated hypothetical documents for dense retrieval
\citep{gao-etal-2023-precise}. Query2doc appends a generated pseudo-document to
the query \citep{wang-etal-2023-query2doc}. MuGI generates multiple references,
pools their contextualized dense representations, and uses query-weighted
concatenation for sparse retrieval \citep{zhang-etal-2024-exploring-best}.
We evaluate all three reproduced pipelines on the same seven datasets as DESA.

We further compare with the rank-based QuDAR-simple variant, which uses the
same generated references and fuses the top-1000 original-sparse (OS),
expanded-sparse (ES), original-dense (OD), and expanded-dense (ED) rankings
with RRF ($k=60$) \citep{kim-etal-2026-qudar}.

\subsection{Mechanism and Robustness Analysis}

We isolate the two channel operators with matched controls. For dense
retrieval, the original sparse ranking is paired with either unprojected
contextual expansion or orthogonal residual expansion. For sparse retrieval,
the original dense ranking is paired with either direct rewriting or
score-product anchoring. A $2\times2$ comparison further evaluates Original,
each proposed operator alone, and their combination.

We then examine how each operator changes its retrieval channel. Dense behavior
is measured by the angle between the original and expanded query vectors.
Sparse behavior is measured by top-20 turnover and by whether relevant
documents move upward in the ranking. We compare these quantities with the
corresponding changes in retrieval effectiveness and access depth using
within-dataset quartiles and Spearman correlations.

Additional experiments vary the number of reference passages over
$\{1,3,5\}$. Score-product anchoring is compared with a binary mask that
retains expanded-query scores only for documents matched by the original
query, separating support preservation from the use of graded original scores.
We also test sparse retrieval using only the generated passages. Robustness
experiments vary the RRF constant over $\{2,20,60,100\}$, replace the generator
and dense encoder, and evaluate the nested TREC-COVID corpora.

\subsection{Statistical Analysis}

The three generation draws are averaged within each query before aggregation.
We report results for each dataset and equal-dataset macro averages. The 95\%
confidence intervals are obtained from 10,000 stratified bootstrap samples,
with queries resampled within each dataset.

The primary comparisons are DESA against Original and DESA
against Shared expansion. We use 10,000 stratified paired sign-flip
randomizations and apply Holm correction across nDCG@10, Recall@20, $L_D$, and
$L_S$ separately for each comparator. The operator controls and the matched
QuDAR comparison use the same draw-averaged, query-level paired procedure, with
Holm correction across nDCG@10 and Recall@20.
